Structured pruning compresses reasoning Large Language Models (LLMs) by removing redundant components from their feed-forward sub-layers, but the pruned model requires recovery training---typically Supervised Fine-Tuning (SFT) or Reinforcement Learning with Verifiable Rewards (RLVR)---to restore reasoning capabilities.
RLVR is widely assumed to be the superior recovery paradigm, yet our results suggest this advantage is attributable to the pruning criterion rather than an inherent property of RL\@.
We propose Recovery-Aware Structured Pruning (\rasp{}), a parameterized scoring function that balances gradient importance with representational diversity measured via Centered Kernel Alignment (CKA).
Across four pruning criteria applied to a 7B reasoning model, we identify a \emph{recovery capacity spectrum}: diversity-preserving criteria sustain active RL training, while importance-dominated criteria exhibit reward collapse.
Despite this stark divergence in training dynamics, SFT and GRPO yield near-identical downstream accuracy at moderate sparsity---a phenomenon we term \emph{reward-accuracy decoupling}.
These findings hold across three benchmarks, with preliminary cross-scale and cross-model evidence.\footnote{Code and data available at \url{https://github.com/anonymous-nlp-2026/recovery_aware_struct_prune}.}
